% Транскрипция семи рукописных листов: ../1a-graphs-handwritten.pdf, стр. 1–7.
% Исходные фотографии 184034–184040 объединены в PDF без перекодирования.
% Уточнено автором: инициалы Г. К.; обозначение набора образующих A(g).
% Сохранены паспорт (2,2,4), обозначения и рассуждения рукописи.
% Зачёркнутые фрагменты и служебные стрелки переходов между листами опущены.
% Сборка из корня репозитория: delivery/_src/render.sh
\documentclass[12pt,a4paper]{article}
\usepackage[margin=20mm]{geometry}
\usepackage{fontspec}
\setmainfont{Times New Roman}
\usepackage{polyglossia}
\setdefaultlanguage{russian}
\usepackage{amsmath,amssymb,array,booktabs}
\usepackage{tikz}
\usetikzlibrary{arrows.meta}
\usepackage{fancyhdr}
\pagestyle{fancy}
\fancyhf{}
\lhead{ДЗ 1a}
\rhead{Гасаев Г. К. \quad КМБО-13-24}
\cfoot{\thepage}
\setlength{\headheight}{15pt}
\setlength{\parindent}{0pt}
\setlength{\parskip}{6pt}
\newcommand{\Aut}{\operatorname{Aut}}
\newcommand{\col}[1]{\begin{pmatrix}#1\end{pmatrix}}
\tikzset{vertex/.style={circle,draw,minimum size=7mm,inner sep=0pt},
  every loop/.style={looseness=8,min distance=13mm}}

\begin{document}
\begin{center}
  {\Large\bfseries Домашнее задание 1a}\\[3pt]
  Гасаев Г. К., КМБО-13-24\\
  16.11.2026
\end{center}

\section*{Условие}
Сначала необходимо сформировать последовательность $a_i$, в которой будут
встречаться только числа $1$, $2$, $3$ и $4$, по следующему правилу.
Надо взять первые несколько букв своей фамилии, добавив в случае слишком
короткой фамилии ещё и имя, заменить каждую букву числом, на единицу большим
остатка от деления на $4$ её номера в алфавите. Нужно выбрать число $k$
таким образом, чтобы было
\[
  \sum_{i=1}^{k}a_i\leq 8,
  \qquad \text{но}\qquad
  \sum_{i=1}^{k+1}a_i>8.
\]
Если число оказалось меньше $8$, надо увеличить одно из наименьших чисел
$a_i$ ($1\leq i\leq k$) на $8-\sum_{i=1}^{k}a_i$.
Если так получилось, что число $k>5$, то надо увеличить наименьшее из чисел
$a_1,\ldots,a_6$ на $10-\sum_{i=1}^{6}a_i$ и затем положить $k=6$.

После того как число $k$ и последовательность чисел $a_1,\ldots,a_k$
определены, нарисуйте все попарно неизоморфные графы (в том числе и несвязные)
с $k$ вершинами, валентности которых равны $a_1,\ldots,a_k$.
Докажите полноту этого списка, вычислив
\[
  \sum_{\Gamma}\frac{1}{|\Aut\Gamma|}
\]
(суммирование по всем нарисованным графам $\Gamma$).

\section*{1. Подготовка вектора $a$}
Фамилия и имя: \textbf{ГАСАЕВГАМИД}.

\begin{center}
\begin{tabular}{cccc}
\toprule
Буква & $N_{\text{алф}}$ & $N_{\text{алф}}\bmod 4$ & $a_i$\\
\midrule
Г & 4 & 0 & 1\\ А & 1 & 1 & 2\\ С & 19 & 3 & 4\\
Е & 6 & 2 & 3\\ В & 3 & 3 & 4\\ М & 14 & 2 & 3\\
И & 10 & 2 & 3\\ Д & 5 & 1 & 2\\
\bottomrule
\end{tabular}
\qquad
\begin{tabular}{c}
Подставляя вместо букв $a_i$, получаем:\\[6pt]
\begin{tabular}{c|ccccccccccc}
 & Г & А & С & А & Е & В & Г & А & М & И & Д\\
\hline
$a_i$ & 1 & 2 & 4 & 2 & 3 & 4 & 1 & 2 & 3 & 3 & 2
\end{tabular}
\end{tabular}
\end{center}
Изначально $a=(1,2,4)$, $k=3$, так как
\[
  1+2+4=7\leq 8,\qquad 1+2+4+2=9>8.
\]
Добавляем $8-\sum_{i=1}^{3}a_i=1$ к наименьшему элементу (первому):
\[
  \boxed{a=(2,2,4).}
\]

\newpage
\section*{2. Найдём все графы с паспортом $\{2,2,4\}$}
Здесь фигурные скобки обозначают мультимножество --- множество,
допускающее повторения.

\subsection*{2.1. Матрица смежности и СЛАУ}
Для начала рассмотрим графы с порядком вершин согласно порядку валентностей
$2,2,4$. Порядок на вершинах обозначаем через имена: $a<b<c$.

Построим матрицу смежности (используем нижний угол):
\[
\begin{array}{c|ccc}
 & a & b & c\\\hline
a & x_1 & & \\
b & x_2 & x_3 & \\
c & x_4 & x_5 & x_6
\end{array}
\]
Тогда паспорт (и выбранный порядок) накладывает такие ограничения на матрицу:
\[
\begin{cases}
2x_1+x_2+x_4=2,\\
x_2+2x_3+x_5=2,\\
x_4+x_5+2x_6=4.
\end{cases}
\]
На диагонали перед каждым коэффициент $2$: это петли, они считаются дважды
в валентностях.

В итоге вектор $X=(x_1,x_2,x_3,x_4,x_5,x_6)^{\mathsf T}$ --- матрица
смежности, свёрнутая в вектор, то есть однозначно описывает какой-то
конкретный граф (а также естественным образом несёт в себе порядок вершин).

Имеем матрицу:
\[
\left(\begin{array}{rrrrrr|r}
2&1&0&1&0&0&2\\
0&1&2&0&1&0&2\\
0&0&0&1&1&2&4
\end{array}\right).
\]
Опустим решение матрицы, способ стандартный для неоднородной СЛАУ:
\[
X=\col{x_1\\x_2\\x_3\\x_4\\x_5\\x_6}
=\col{0\\0\\0\\2\\2\\0}
+\alpha\col{1\\-1\\0\\-1\\1\\0}
+\beta\col{0\\-1\\1\\1\\-1\\0}
+\gamma\col{0\\1\\0\\-1\\-1\\1}.
\]
Добавляя к $X$ условия на целочисленность и неотрицательность всех $x_i$
(это валентности):
\[
x_i\geq0\quad\Longleftrightarrow\quad
\begin{cases}
\alpha\geq0,\\
-\alpha-\beta+\gamma\geq0,\\
\beta\geq0,\\
2-\alpha+\beta-\gamma\geq0,\\
2+\alpha-\beta-\gamma\geq0,\\
\gamma\geq0,\\
\alpha,\beta,\gamma\in\mathbb Z
\end{cases}
\quad\Longleftrightarrow\quad
\begin{cases}
\alpha\geq0,\quad\beta\geq0,\\
\alpha+\beta\leq\gamma,\\
\gamma\leq2-\alpha+\beta,\\
\gamma\leq2+\alpha-\beta,\\
\alpha,\beta,\gamma\in\mathbb Z.
\end{cases}
\]

\newpage
\subsection*{2.2. Перебор решений}
\[
\gamma\leq\min(2-\alpha+\beta,2+\alpha-\beta)=2-|\alpha-\beta|.
\]
Поэтому система равносильна
\[
\begin{cases}
\alpha\geq0,\quad\beta\geq0,\\
\alpha+\beta\leq\gamma\leq2-|\alpha-\beta|,\\
\alpha,\beta,\gamma\in\mathbb Z
\end{cases}
\quad\Longleftrightarrow\quad
\begin{cases}
\alpha\geq0,\quad\beta\geq0,\\
\gamma\in[\alpha+\beta,\,2-|\alpha-\beta|],\\
\alpha+\beta+|\alpha-\beta|\leq2,\\
\alpha,\beta,\gamma\in\mathbb Z.
\end{cases}
\]
Так как $\alpha+\beta+|\alpha-\beta|=2\max(\alpha,\beta)$, получаем
\[
\begin{cases}
\alpha,\beta\in\{0,1\},\\
\gamma\in[\alpha+\beta,\,2-|\alpha-\beta|],\\
\gamma\in\mathbb Z.
\end{cases}
\]
Перебираем все $(\alpha,\beta)$:
\begin{center}
\begin{tabular}{ccc}
\toprule
$(\alpha,\beta)$ & Ограничение на $\gamma$ & Целые $\gamma$\\
\midrule
$(0,0)$ & $[0,2]$ & $0,1,2$\\
$(0,1)$ & $[1,1]$ & $1$\\
$(1,0)$ & $[1,1]$ & $1$\\
$(1,1)$ & $[2,2]$ & $2$\\
\bottomrule
\end{tabular}
\end{center}
Таким образом,
\[
(\alpha,\beta,\gamma)\in
\{(0,0,0),(0,0,1),(0,0,2),(0,1,1),(1,0,1),(1,1,2)\}.
\]
Подставляя в $X$, получаем:
\[
X\in\left\{
\col{0\\0\\0\\2\\2\\0},
\col{0\\1\\0\\1\\1\\1},
\col{0\\2\\0\\0\\0\\2},
\col{0\\0\\1\\2\\0\\1},
\col{1\\0\\0\\0\\2\\1},
\col{1\\0\\1\\0\\0\\2}
\right\}.
\]
Получили шесть графов, подходящих под паспорт с порядком $(2,2,4)$.
Но мы не учли порядки $(2,4,2)$ и $(4,2,2)$.
Они получаются из найденных графов перестановками вершин.
Все новые векторы будут изоморфны каким-то из найденных, однако и среди
найденных могут быть изоморфные. То есть система полна, но не минимальна
(точнее, мы пока не знаем, минимальна ли).

\newpage
\section*{3. Классы изоморфизма}
Любые перестановки вершин (их в данном случае $3!=6$, для трёх вершин)
сохраняют структуру рёбер (и паспорт), то есть ведут к изоморфному графу.
И каждой из них соответствует перестановка координат $x_i$ матрицы
(вектора) смежности:
\begin{center}
\begin{tabular}{ccc}
\toprule
Перестановка вершин & Соответствующая перестановка $x_i$ & Вектор после перестановки\\
\midrule
$(ab)$ & $(13)(45)$ & $(x_3,x_2,x_1,x_5,x_4,x_6)$\\
$(ac)$ & $(16)(25)$ & $(x_6,x_5,x_3,x_4,x_2,x_1)$\\
$(bc)$ & $(24)(36)$ & $(x_1,x_4,x_6,x_2,x_5,x_3)$\\
$(abc)$ & $(136)(254)$ & $(x_6,x_4,x_1,x_5,x_2,x_3)$\\
$(acb)$ & $(163)(245)$ & $(x_3,x_5,x_6,x_2,x_4,x_1)$\\
\bottomrule
\end{tabular}
\end{center}
Теперь по найденным из СЛАУ шести векторам и перестановкам на $x_i$,
сохраняющим изоморфизм, восстановим все графы паспорта $\{2,2,4\}$
и их связи друг с другом.

В таблицах первая строка задаёт исходный вектор, остальные --- его образы.
Например, запись $000220$ означает $(0,0,0,2,2,0)$.

\begin{center}
\begin{minipage}[t]{.47\textwidth}
\centering\textbf{Класс 1}\\[5pt]
\begin{tabular}{cc}
Перестановки & Вектор\\\hline
$\mathrm{id},(ab)$ & $000220$\\
$(bc),(acb)$ & $020020$\\
$(abc),(ac)$ & $020200$
\end{tabular}

\medskip
Представитель: $000220$.

\begin{tikzpicture}
\node[vertex] (a) at (0,1.5) {$a$};
\node[vertex] (b) at (2,1.5) {$b$};
\node[vertex] (c) at (1,0) {$c$};
\draw (a) to[bend left=18] (c) (a) to[bend right=18] (c);
\draw (b) to[bend left=18] (c) (b) to[bend right=18] (c);
\end{tikzpicture}
\end{minipage}\hfill
\begin{minipage}[t]{.47\textwidth}
\centering\textbf{Класс 2}\\[5pt]
\begin{tabular}{cc}
Перестановки & Вектор\\\hline
$\mathrm{id},(ab)$ & $010111$\\
$(bc),(acb)$ & $011110$\\
$(abc),(ac)$ & $110110$
\end{tabular}

\medskip
Представитель: $010111$.

\begin{tikzpicture}
\node[vertex] (a) at (0,1.5) {$a$};
\node[vertex] (b) at (2,1.5) {$b$};
\node[vertex] (c) at (1,0) {$c$};
\draw (a)--(b)--(c)--(a) (c) edge[loop below] (c);
\end{tikzpicture}
\end{minipage}
\end{center}

\newpage
\begin{center}
\begin{minipage}[t]{.47\textwidth}
\centering\textbf{Класс 3}\\[5pt]
\begin{tabular}{cc}
Перестановки & Вектор\\\hline
$\mathrm{id},(ab)$ & $020002$\\
$(bc),(acb)$ & $002200$\\
$(abc),(ac)$ & $200020$
\end{tabular}

\medskip
Представитель: $020002$.

\begin{tikzpicture}
\node[vertex] (a) at (0,1.8) {$a$};
\node[vertex] (b) at (2,1.8) {$b$};
\node[vertex] (c) at (1,0) {$c$};
\draw (a) to[bend left=18] (b) (a) to[bend right=18] (b);
\draw (c) edge[loop left] (c) (c) edge[loop right] (c);
\end{tikzpicture}
\end{minipage}\hfill
\begin{minipage}[t]{.47\textwidth}
\centering\textbf{Класс 4}\\[5pt]
\begin{tabular}{cc}
Перестановки & Вектор\\\hline
$\mathrm{id},(ab)$ & $101002$\\
$(bc),(acb)$ & $102001$\\
$(ac),(abc)$ & $201001$
\end{tabular}

\medskip
Представитель: $101002$.

\begin{tikzpicture}
\node[vertex] (a) at (0,1.8) {$a$};
\node[vertex] (b) at (2,1.8) {$b$};
\node[vertex] (c) at (1,0) {$c$};
\draw (a) edge[loop above] (a) (b) edge[loop above] (b);
\draw (c) edge[loop left] (c) (c) edge[loop right] (c);
\end{tikzpicture}
\end{minipage}
\end{center}

\subsection*{Класс 5}
\begin{center}
\begin{minipage}{.46\textwidth}
\centering
\begin{tabular}{cc}
Перестановка & Вектор\\\hline
$\mathrm{id}$ & $100021$\\
$(ab)$ & $001201$\\
$(acb)$ & $021001$\\
$(ac)$ & $120001$\\
$(abc)$ & $101200$\\
$(bc)$ & $101020$
\end{tabular}
\end{minipage}\hfill
\begin{minipage}{.46\textwidth}
\centering
Представитель: $100021$.

\begin{tikzpicture}
\node[vertex] (a) at (0,1.8) {$a$};
\node[vertex] (b) at (2,1.8) {$b$};
\node[vertex] (c) at (1,0) {$c$};
\draw (a) edge[loop above] (a) (c) edge[loop below] (c);
\draw (b) to[bend left=18] (c) (b) to[bend right=18] (c);
\end{tikzpicture}
\end{minipage}
\end{center}
Как видно, в найденных из СЛАУ решениях было два изоморфных графа:
$100021$ и $001201$ из класса 5.

Таким образом, система представителей каждого класса
\[
\Gamma=\left\{
\col{0\\0\\0\\2\\2\\0},
\col{0\\1\\0\\1\\1\\1},
\col{0\\2\\0\\0\\0\\2},
\col{1\\0\\1\\0\\0\\2},
\col{1\\0\\0\\0\\2\\1}
\right\}
\]
полна и минимальна.

\newpage
\section*{4. Проверка полноты через массу системы}
Теперь проверим это, вычислив массу полной системы по формуле (пункт 4.2)
\[
M_{\text{полн}}=\frac{PQ}{(2N)!},\qquad
N=\frac12\sum_i n_i,\qquad
P=\frac{(2N)!}{N!\,2^N},
\]
\[
Q=\frac{(2N)!}{(n_1!)^{m_1}m_1!\cdots(n_s!)^{m_s}m_s!}.
\]
Здесь $n_1,\ldots,n_s$ --- встречающиеся валентности,
$m_1,\ldots,m_s$ --- частоты, с которыми они встречаются в векторе паспорта.
Сравним с массой нашей системы, вычисленной по определению (пункт 4.1):
\[
M_{\Gamma}=\sum_{g\in\Gamma}\frac{1}{|\Aut g|}.
\]

\subsection*{4.1. Автоморфизмы графов из $\Gamma$}
Обозначения:
\begin{itemize}
\item $T$ --- обмен вершин $a,b$, переставляющий рёбра и петли,
сохраняя их номера;
\item $P_{uv}$ --- обмен параллельных рёбер между $u$ и $v$;
\item $P_c$ --- обмен двух петель у $c$;
\item $F_u$ --- переворот петли у $u$ ($F_{u,1}$, $F_{u,2}$, если две петли).
\end{itemize}

\begin{center}
\renewcommand{\arraystretch}{1.4}
\begin{tabular}{clc}
\toprule
Граф $g$ & Основание $A(g)$ & Количество автоморфизмов $2^{|A(g)|}=|\Aut g|$\\
\midrule
$000220$ & $P_{ac},P_{bc},T$ & $2^3=8$\\
$010111$ & $F_c,T$ & $2^2=4$\\
$020002$ & $P_{ab},F_{c,1},F_{c,2},P_c,T$ & $2^5=32$\\
$101002$ & $F_a,F_b,F_{c,1},F_{c,2},P_c,T$ & $2^6=64$\\
$100021$ & $P_{bc},F_a,F_c$ & $2^3=8$\\
\bottomrule
\end{tabular}
\end{center}
Тогда
\[
M_{\Gamma}=\sum_{g\in\Gamma}\frac{1}{|\Aut g|}
=\frac18+\frac14+\frac1{32}+\frac1{64}+\frac18
=\boxed{\frac{35}{64}}.
\]

\newpage
\subsection*{4.2. Масса полной системы по комбинаторной формуле}
\[
(n_1,n_2,n_3)=(2,2,4),
\]
\[
N=\frac12\sum_i n_i=\frac{2+2+4}{2}=4,
\qquad 2N=8.
\]
\[
P=\frac{(2N)!}{N!\,2^N}
=\frac{8!}{4!\,2^4}=105.
\]
\[
Q=\frac{8!}{(2!)^2\,4!\,2!}=210.
\]
\[
M_{\text{полн}}=\frac{PQ}{(2N)!}
=\frac{105\cdot210}{8!}
=\boxed{\frac{35}{64}}.
\]

\section*{5. Вывод}
Как видно, $M_{\text{полн}}=M_{\Gamma}$, что ещё раз доказывает,
что $\Gamma$ --- полная система.
\end{document}
